\section{Semiring contraction}

We here introduce with commutative semirings the algebraic structure required for the main separability property of contractions, and therefore for message-passing algorithms.

See \cite{liu_computing_2023,liu_tropical_2021} for applications and \cite{gondran_graphs_2008} for algebraic and graph-theoretic foundations.

\subsection{Foundations}

\begin{definition}
    A commutative semiring is a tuple $\gensemiring$, where $S$ is a set, $0,1\in S$, and
    \begin{align*}
        \srplus \defcols S \times S \rightarrow S \quad, \quad \srmul \defcols S \times S \rightarrow S \, ,
    \end{align*}
    such that the following properties hold:
    \begin{itemize}
        \item[(i)] $(S,\srplus,0)$ is a commutative monoid\footnote{That is associativity, commutativity of $\srplus$ and identity of $s\srplus 0=s$ for all $s\in S$.}.
        \item[(ii)] $(S,\srmul,1)$ is a commutative monoid.
        \item[(iii)] $\srmul$ distributes over $\srplus$, that is for any $s_0,s_1,s_2\in S$
        \begin{align*}
            s_0 \srmul (s_1 \srplus s_2)
            = (s_0 \srmul s_1) \srplus (s_0 \srmul s_2) \, .
        \end{align*}
        \item[(iv)] $0$ annihilates $\srmul$, that is for any $s_0\in S$
        \begin{align*}
            0 \srmul s_0 = 0 \, .
        \end{align*}
    \end{itemize}
\end{definition}

We understand $S$ as the set of coordinates in a tensor.

\begin{definition}
    A tensor with coordinates in $S$ of order $\catorder\in\nn$ and leg dimensions $\catdimof{[\catorder]}=(\catdimof{0},\ldots,\catdimof{\catorder-1})$, where $\catdimof{\catenumerator}\in\nn$ for $\catenumeratorin$, is a map
    \begin{align*}
        \hypercore \defcols \facstates \rightarrow S \, .
    \end{align*}
    A tensor network with coordinates in $S$ is a hypergraph $\graph=(\nodes,\edges)$, where each node is decorated by a dimension $\catdimof{\node}$ and each edge by a tensor of order $\cardof{\edge}$ and leg dimensions $\catdimof{\edge}$.
\end{definition}


%\section{Contractions}

\begin{definition}
    Given a semiring $\gensemiring$ and a tensor network $\extnetasset$ with coordinates in $S$, the contraction of the tensor network keeping the variables to a subset $\arbset\subset\nodes$ open is the tensor
    \begin{align*}
        \srcontractionofwrt{
            \extnetasset
        }{\catvariableof{\arbset}}{\gensemiring}
    \end{align*}
    with coordinates at indices $\catindexof{\arbset}$ by
    \begin{align*}
        \srcontractionofwrt{
            \extnetasset
        }{\catvariableof{\arbset}}{\gensemiring}
        = \bigoplus_{\catindexof{\nodes\backslash\arbset}} \bigotimes_{\edgein} \hypercoreofat{\edge}{\indexedcatvariableof{\edge}} \, .
    \end{align*}
\end{definition}

The distributive property $(iii)$ of the semiring ensures, that we can compute contractions by sub-contractions.

\begin{theorem}
    Given a disjoint partition of $\edges$ into $\edgesof{0},\edgesof{1}$, and a subset $\arbset\subset\nodes$, we define the sets $\nodesof{0}$ and $\nodesof{1}$ as the nodes appearing in $\edgesof{0}$, respectively $\edgesof{1}$.
    We then have
    \begin{align*}
        \srcontractionofwrt{
            \extnetasset
        }{\catvariableof{\arbset}}{\gensemiring}
        =
        \srcontractionofwrt{
            \srcontractionofwrt{
                \{\hypercoreofat{\edge}{\catvariableof{\edge}} \wcols \edge\in\edgesof{0}\}
            }{\catvariableof{\nodesof{0}\cap(\arbset\cup\nodesof{1})}}{\gensemiring},
            \srcontractionofwrt{
                \{\hypercoreofat{\edge}{\catvariableof{\edge}} \wcols \edge\in\edgesof{1}\}
            }{\catvariableof{\nodesof{1}\cap(\arbset\cup\nodesof{0})}}{\gensemiring}
        }{\catvariableof{\arbset}}{\gensemiring} \, .
    \end{align*}
\end{theorem}
\begin{proof}
%        We use the commutative properties to split the summations and the distributive property to move the $\bigoplus$ to closed variables only in $\nodesof{0}$ respectively $\nodesof{1}$ to the first respectively second sub-contraction.
    We use the associativity and commutativity of $\srmul$ to get for any index $\catindexof{\nodes}$
    \begin{align*}
        \bigotimes_{\edgein}\hypercoreofat{\edge}{\indexedcatvariableof{\edge}}
        = \left(\bigotimes_{\edge\in\edgesof{0}}\hypercoreofat{\edge}{\indexedcatvariableof{\edge}}\right)
        \srmul \left(\bigotimes_{\edge\in\edgesof{1}}\hypercoreofat{\edge}{\indexedcatvariableof{\edge}}\right) \, .
    \end{align*}
    Let us now keep $\catindexof{\arbset}$ fixed and investigate the $\srplus$ operation over the indices $\catindexof{\nodes\backslash\arbset}$.
    Distributivity of $\srmul$ over $\srplus$ implies
    \begin{align*}
        \sum_{\catindexof{\nodes\backslash\arbset}} \bigotimes_{\edgein}\hypercoreofat{\edge}{\indexedcatvariableof{\edge}}
        = \sum_{\catindexof{(\nodesof{0}\cap\nodesof{1})\backslash \arbset}}
        \left(
            \sum_{\catindexof{\nodesof{0}\backslash(\arbset\cup\nodesof{1})}}\bigotimes_{\edge\in\edgesof{0}}\hypercoreofat{\edge}{\indexedcatvariableof{\edge}}
        \right) \srmul
        \left(
            \sum_{\catindexof{\nodesof{1}\backslash(\arbset\cup\nodesof{0})}}\bigotimes_{\edge\in\edgesof{1}}\hypercoreofat{\edge}{\indexedcatvariableof{\edge}}
        \right)
    \end{align*}
    which is equivalent to
    \begin{align*}
        \srcontractionofwrt{
            \extnetasset
        }{\indexedcatvariableof{\arbset}}{\gensemiring}
        =
        \srcontractionofwrt{
            \srcontractionofwrt{
                \{\hypercoreofat{\edge}{\catvariableof{\edge}} \wcols \edge\in\edgesof{0}\}
            }{\catvariableof{\nodesof{0}\cap(\arbset\cup\nodesof{1})}}{\gensemiring},
            \srcontractionofwrt{
                \{\hypercoreofat{\edge}{\catvariableof{\edge}} \wcols \edge\in\edgesof{1}\}
            }{\catvariableof{\nodesof{1}\cap(\arbset\cup\nodesof{0})}}{\gensemiring}
        }{\indexedcatvariableof{\arbset}}{\gensemiring} \, .
    \end{align*}
    Since therefore both tensors of the claimed contraction equation coincide on all coordinates, they are equal.
\end{proof}

As a consequence, message passing schemes such as variable elimination generalize to semiring tensor network contractions.

\subsection{Examples}

\subsubsection{Sum-product TN}\label{sec:complexSPContraction}

Standard TN use the semiring
\begin{align*}
(\rr,+,\cdot,0,1)
    \, .
\end{align*}


Variants are:
\begin{itemize}
    \item Restriction of $\rr$ to non-negative numbers $[0,\infty)$: Ensuring the probabilistic interpretation
    \item Generalization of $\rr$ to complex numbers $\mathbb{C}$: Used in quantum computing
\end{itemize}

\subsubsection{Constraint Satisfaction}

We take $S=\{0,1\}$ and interpret $0$ as $\falsesymbol$, $1$ as $\truesymbol$.
We define the disjunction and conjunction operations by
\begin{align*}
    \lor \left[\catvariableof{0},\catvariableof{1}\right]
    = \coloredmatrixof{0 & 1 \\ 1 & 1}{\catvariableof{0},\catvariableof{1}}
\end{align*}
and
\begin{align*}
    \land \left[\catvariableof{0},\catvariableof{1}\right]
    = \coloredmatrixof{0 & 0 \\ 0 & 1}{\catvariableof{0},\catvariableof{1}} \, .
\end{align*}

Tensors over $\{0,1\}$ are interpreted as constraints, and their tensor networks as constraint networks.

The satisfiability of a constraint network $\extnetasset$ is the contraction
\begin{align*}
    \bigvee_{\catindexof{\nodes}} \bigwedge_{\edgein} \hypercoreofat{\edge}{\indexedcatvariableof{\edge}}
    = \srcontractionwrt{
        \extnetasset
    }{\satsemiring} \, .
\end{align*}

As an additional property, we have invariance of adding subcontractions to the network.
In general tensor networks, this holds only for the support of subcontractions (being the operator mapping $0$ to $0$ and any other $s\in S$ to $1$).

\subsubsection{Maximum calculus}

The maximum semiring is the tuple
\begin{align*}
(\rr,\mathrm{max},\cdot,0,1)
    \, .
\end{align*}

The full contraction is the maximal coordinate in the tensor
\begin{align*}
    \srcontractionwrt{\hypercorewith}{(\rr,\mathrm{max},\cdot,0,1)}
    = \max_{\shortcatindicesin} \hypercoreat{\indexedshortcatvariables} \, .
\end{align*}
When leaving variables open in the contraction, we get the maximal coordinate in the specific slices of the tensor, namely
\begin{align*}
    \srcontractionofwrt{\hypercorewith}{\indexedcatvariableof{\secnodes}}{(\rr,\mathrm{max},\cdot,0,1)}
    = \max_{\catindexof{\nodes\backslash\secnodes}} \hypercoreat{\indexedshortcatvariables} \, .
\end{align*}

Analogous properties has the minimal semiring
\begin{align*}
(\rr,\mathrm{min},\cdot,0,1)
    \, .
\end{align*}


These schemes only work, when taking the minimal coordinate keeping all variables open.
When searching for the maximal/minimal coordinate of a contraction with some variables closed, we need to first coarse grain the tensor network.

These contractions are used in encoding-decoding schemes \cite{mezard_information_2009,mackay_information_2003}.
Message passing schemes involving the maximum semiring are known as max-product algorithms \cite{koller_probabilistic_2009}.

\subsubsection{Tropical semiring}

The tropical semiring is
\begin{align*}
    \mathbb{T}
    = (\rr\cup\{\infty\},\mathrm{min},+,\infty,0) \, .
\end{align*}

The contraction of a tensor gives the minimum
\begin{align*}
    \srcontractionwrt{\hypercorewith}{\mathbb{T}}
    = \min_{\shortcatindicesin} \hypercoreat{\shortcatvariables} \, .
\end{align*}

The tensor product in the tropical semiring sums the leg vector coordinates, that is
\begin{align*}
    \srcontractionofwrt{
        \{\legcoreofat{\catenumerator}{\catvariableof{\catenumerator}} \wcols \catenumeratorin\}
    }{\indexedshortcatvariables}{\mathbb{T}}
    = \sum_{\catenumeratorin} \legcoreofat{\catenumerator}{\indexedcatvariableof{\catenumerator}} \, .
\end{align*}

Further application in \cite{hu_tropical_2026} and on spin glasses \cite{liu_tropical_2021}.

\subsubsection{Galois field of two elements}\label{sec:semiringGalois}

The Galois field of two elements is the semiring
\begin{align*}
    \galoissemiring
    \coloneqq (\ozset,\oplus,\cdot,0,1) \, ,
\end{align*}
where $\oplus$ denotes the sum $\modspace 2$ and $\cdot$ the standard multiplication.

\begin{example}{Contractions in the semiring $\galoissemiring$}
    While $\contraction{\formulawith}$ is the model count, the contraction with respect to $\galoissemiring$ decides whether the number of models is odd or even:
    \begin{align*}
        \srcontractionwrt{\formulawith}{\galoissemiring}
        &= \bigoplus_{\shortcatindicesin} \formulaat{\indexedshortcatvariables} \\
        &= \cardof{\{\shortcatindices \wcols \formulaat{\indexedshortcatvariables}=1\}} \modspace 2 \, .
    \end{align*}
\end{example}


We will investigate that semiring further in \secref{sec:reedMullerExpansion}, where we relate $\cpformat$ decompositions with respect to that semiring to Reed-Muller expansions.

\subsection{Tensor network formats}

Based on the semiring contraction, we can define tensor network formats for arbitrary semirings.
\begin{itemize}
    \item $\elformat$ depends only on the multiplicative structure of the semiring. Therefore equivalent for all the semirings with the standard product.
    \item $\cpformat$ format is the semiring sum of $\elformat$. Examples in the galois field semiring in \secref{sec:reedMullerExpansion}.
\end{itemize}